\section{Proofs for the paid-support policy safeguard}
\label{supp:supportaction122}

All claims below are deterministic statements about the original economic
operator. Conditional sampling and population inference, if used in an
experiment, require a separately frozen protocol. None of the proofs changes
the current information filtration or the terminal controller.

\subsection{Original-gradient cap and the producer's support semantics}

The original state Lipschitz bound \eqref{eq:lipschitz95}, specialized to
$d=2$, and the action exposure $\boldsymbol b=(1/2,1/4)$ give
\[
 |V_1(F(x,a,z))-V_1(F(x,b,z))|
 \le \frac{15}{2}\frac34|a-b|=\frac{45}{8}|a-b|.
\]
Monotonicity and cash invariance imply the sup-norm Lipschitz bound
$|\rho(U)-\rho(V)|\le\|U-V\|_\infty$ for the original AVaR.
Since $0\le h'(a)=2a+16a^3\le3/4$ on $[0,1/4]$, it follows that
\[
 |Q_x(a)-Q_x(b)|\le
 \left(\frac34+\frac{15}{16}\frac{45}{8}\right)|a-b|
 =\frac{771}{128}|a-b|<7|a-b|.
\]
The established differentiability of the original $Q_x$ therefore gives
$|Q_x'|\le771/128<7$, including inward endpoint derivatives.
This argument holds for both declared risk levels and every original state;
it does not estimate a derivative cap from sample gradients.

Here is why the existing original-risk payload supplies
\eqref{eq:paidinformation122}. Freeze any feasible risk density constant
on the original innovation cells $B_i$, with exact masses $w_i$ and
midpoints $z_i$. At $y_i=F(x,a,z_i)$, retain a global affine terminal
minorant $V_1(y)\ge\ell_i+p_i^{\mathsf T}(y-y_i)$.
Within a cell, the original transition satisfies
\[
 F(x,b,z)-F(x,a,z_i)
      =\boldsymbol b(b-a)+\boldsymbol s(z-z_i).
\]
The integral of $z-z_i$ under its constant density is zero.
Thus the same fixed feasible risk witness, evaluated at an arbitrary
feasible action $b$, gives
\[
 Q_x(b)\ge S(x)+h(b)+
  \beta\sum_iw_i\{\ell_i+p_i^{\mathsf T}\boldsymbol b(b-a)\}.
\]
Using $h(b)\ge h(a)+h'(a)(b-a)+(b-a)^2$ proves the global quadratic
support with
\[
 W=S(x)+h(a)+\beta\sum_iw_i\ell_i,\qquad
 p=h'(a)+\beta\sum_iw_i\boldsymbol b^{\mathsf T}p_i.
\]
An outward lower $W$ and an interval enclosing the actual $p$ preserve it.
The witness may have been selected by sorting lower midpoint values at $a$;
its weights are held fixed in the displayed comparison. It is not
implicitly differentiated as an optimizing density at another action.
The same proof covers a feasible active-tail density.

The original upper $U$ comes from the convex-order endpoint-mixture
calculation for $V_1$ under the original continuous law. The implementation
exposes its separately recomputed upper
\texttt{(base + BETA * upper\_risk).hi} before the terminal-controller
allowance. It preserves the actual lower witness separately from the
reported policy-interval lower endpoint; they need not have identical
rounding. This avoids treating a generic policy upper bound as an exact
optimal-future value or subtracting an unidentified error term.

\begin{proof}[Proof of Theorem~\ref{thm:supportgradient122}]
For $0<h\le A-a$, derivative Lipschitzness and the lower support give
\[
 W+ph+\frac\mu2h^2
 \le q(a+h)
 \le q(a)+q'(a)h+\frac L2h^2
 \le U+q'(a)h+\frac L2h^2.
\]
Rearranging gives $q'(a)\ge p-e/h-ch$ and therefore the lower
inequality after $p\ge p_-$ and minimization in $h$.
For $0<h\le a$ use $a-h$ in the same calculation. Dividing by the
negative displacement gives $q'(a)\le p+e/h+ch\le p_++e/h+ch$.
This proves the upper inequality. Integrating the continuous derivative
proves the smooth upper Taylor inequality used here; strong convexity
supplies the lower quadratic term. One-sided boundary derivatives have
the indicated inward interpretation.

For $e,c>0$, $e/h+ch$ decreases until $h=\sqrt{e/c}$ and increases
thereafter, which proves \eqref{eq:psi122}. At $e=0$ its infimum on
$(0,r]$ is zero. If $c=0$, the infimum is $e/r$, including zero.
When $r=0$ there is no admissible displacement in that derivation and
no finite inference on that side.

We next show sharpness, which concerns the stated information rather than
the particular investment formula. Keep an exact slope $p$, the support
level $W$, the upper level $U$, and the action interval fixed. Consider
\[
 q_g(a+d)=U+gd+\frac L2d^2.
\]
This function has $L$-Lipschitz derivative and is $\mu$-strongly convex
because $L\ge\mu$. Its difference from the supplied support is
\[
 q_g(a+d)-\{W+pd+\tfrac\mu2d^2\}
       =e+(g-p)d+cd^2.
\]
For all positive feasible $d$, nonnegativity of this expression is
exactly $g-p\ge-\inf_{0<d\le A-a}(e/d+cd)$.
For all negative feasible $d$ it is exactly
$g-p\le\inf_{0<-d\le a}(e/(-d)+c(-d))$.
At $d=0$ it is $e\ge0$. Thus every derivative in
\[
 [p-\Psi(e,c,A-a),\ p+\Psi(e,c,a)]
\]
is realized by this quadratic with the same supplied information,
and neither side can be tightened. At $a=0$, arbitrarily large positive
$g$ remain feasible, and at $a=A$, arbitrarily negative $g$ do so.
If an extra original-gradient bound is intersected, validity persists;
sharpness after additional model-specific restrictions is not asserted.

For a concrete boundary warning, set $A=1/4$, $a=0$, $\mu=2$,
$L=10$, $W=U=0$, and $p=-1$. Both
$q_-(b)=-b+5b^2$ and $q_+(b)=b+5b^2$ lie above the common support
$-b+b^2$ on $[0,1/4]$. Their inward derivatives at zero have opposite
signs. Treating the touching boundary support slope as the actual
derivative would falsely certify descent for $q_+$.
\end{proof}

\begin{proof}[Proof of Theorem~\ref{thm:payloadsafeguard122}]
Integration of the derivative inequalities, or the defining strong
convexity and smoothness inequalities, gives for every feasible $d$
\[
 q'(a)d+\tfrac\mu2d^2
 \le q(a+d)-q(a)
 \le q'(a)d+\tfrac L2d^2.
\]
Taking the extrema over the valid derivative interval proves
$\underline m(d)\le q(a+d)-q(a)\le\overline m(d)$.
The future-policy difference belongs to $[-\delta,\delta]$,
so addition proves \eqref{eq:actualcontrast122}. This step uses no
convexity, continuity, or differentiability of the implemented future
policy's error.

The saved upper value gives
$J_x(a+d)\le U+\overline m(d)+\delta=U_b$.
The quadratic lower support and $J_x\ge q$ give the lower endpoint
$L_b$. For every feasible action, the saved support minorizes $q$.
Its infimum for the actual unknown slope is at least the smaller of
the infima for $p_-$ and $p_+$, since for each displacement the affine
function of $p$ is bounded below by its endpoint minimum. Thus
$B\le\min_{b\in[0,A]}q(b)=V_0(x)$. Taking the maximum with another
valid original-optimum lower bound preserves this inequality.
Subtracting proves the regret bound. A feasible implemented policy
cannot beat $V_0$, by the monotone original recursion and its feasible
last-date controller, which gives the nonnegative side.

On $d\ge0$, the robust upper model equals
$g_+d+Ld^2/2$; on $d\le0$ it equals
$g_-d+Ld^2/2$. Each is convex. If the interval straddles zero, both
pieces are nonnegative and $d=0$ minimizes them. If $g_->0$, the
nonpositive part has its unconstrained minimizer $-g_-/L$, clipped
to $[-a,0]$; the positive part cannot improve on zero.
The negative-gradient case is symmetric. This proves the displayed
proposal formula. The three actual acceptance checks concern only the
realized feasible action, so a rounded action is valid when its newly
computed inequalities pass. No bound for a different action is reused.

A strict negative upper contrast proves actual improvement. The two
absolute tests prove qualification. At an accepted sequence, sum the
statewise inequalities $J(a_{j-1})-J(a_j)\ge
-\overline m_j(d_j)-\delta$; the actual costs telescope. Retention
steps add zero. Every accepted action is feasible and qualified by its
own three tests, completing the induction. A finite measurable rule
with fixed tie and failure conventions supplies an economic policy.
The test is valid even when the proposal was produced by a nonconvex
critic, since the premises concern the original $q$.

For the claimed information sharpness only, take an interior anchor,
an exact supplied support slope, and the exact untruncated interval
\eqref{eq:gradientinterval122}, before imposing any further original-model
restriction. Use the quadratic $q_g$ in the preceding proof with the endpoint
of that exact feasible gradient interval maximizing $gd$. It gives
$q_g(a+d)-q_g(a)=\overline m(d)$.
There is no representation-dependent slack to remove under those data.
The additional future-policy floor can also be attained within the
interface $0\le\eta\le\delta$ at two distinct actions by setting
$\eta(a)=0,\eta(a+d)=\delta$. This reproduces the retained R120
future-information limitation and is not a new claim about what the
particular terminal controller must realize.
\end{proof}

\begin{proposition}[A checkable interior precision ladder]
\label{prop:precisionladder122}
Let $a\in(0,A)$, $L>\mu$, and suppose level $k$ supplies the information
in \eqref{eq:paidinformation122}, with $e_k\le e_0 4^{-k}$ and
$p_{k,+}-p_{k,-}\le w_0 2^{-k}$. At every level satisfying
$e_0 4^{-k}\le (L-\mu)\min\{a,A-a\}^2/2$, the exact derivative
interval has width at most $D2^{-k}$, where
$D=w_0+2\sqrt{2(L-\mu)e_0}$.
Let $\Gamma_k=-\min_d\overline m_k(d)$ be the unrounded robust
smooth-model gain. Then
\[
 0\le G_L(a,q'(a))-\Gamma_k\le AD2^{-k}.
\]
With a realized rounding displacement error at most $\rho$, an additional
conservative loss allowance $(G+LA)\rho$ is valid whenever the gradient
interval used in the robust model lies in $[-G,G]$.
\end{proposition}

\begin{proof}
The interior condition permits the minimizing positive $h$ on both
sides. The uncertainty on each side is
$\sqrt{2(L-\mu)e_k}$, which yields the width bound.
Let $I=[l,u]$ be that interval and let $g=q'(a)\in I$.
Here $I$ is the exact-infimum interval; a larger finite-radius numerical
interval requires its actual width in this argument. The scalar $g_0$ in
$I$ closest to zero is $l$ if $l>0$, $u$ if
$u<0$, and zero otherwise. The formula for the robust proposal shows
$\Gamma_k=G_L(a,g_0)$. One may also see this by separating the
positive and negative displacement domains.
If $g_0=0$, the robust gain is zero; otherwise it uses the same
sign as every $g\in I$ and the smallest magnitude.
Therefore $G_L(a,g_0)\le G_L(a,g)$.
The maximum of the functions $-gd-Ld^2/2$ is $A$-Lipschitz in $g$,
since every feasible $d$ satisfies $|d|\le A$.
Also $|g-g_0|\le u-l\le D2^{-k}$. These prove the assertion.

For rounding, on the feasible displacement interval each piece of
$\overline m$ has derivative of magnitude at most $G+LA$.
The pieces meet continuously at zero, so the same bound is global.
A displacement perturbation of at most $\rho$ therefore increases
the robust upper model by at most $(G+LA)\rho$.
In actual code it is preferable to reevaluate the outward model at
the realized action, rather than to spend this extra allowance.
\end{proof}

This proposition gives a finite gain stopping bound under explicit
oracle guarantees and a strict margin. Its scope must be kept separate
from the qualification-width test: as $e$ and slope uncertainty tend
to zero, \eqref{eq:candidateinterval122} can still have width
$(L-\mu)d^2/2+\delta$. Additional candidate information is necessary
when that curvature width exceeds the prescribed requirement.
The $L=\mu$ case has $\Psi(e,0,r)=e/r$ and admits the corresponding
faster algebraic bound, but no nonzero-capacity or interior room can
be inferred at a boundary.

\begin{proposition}[A priced stopping region]
\label{prop:pricedstop122}
Suppose an implemented incumbent has certified original regret at most
$K\ge0$. Suppose every admissible continuation of a single-decision
solver incurs at least $c_*$ extra computation units before it can
change the decision, these units cost $\lambda\ge0$ per unit in
original-objective units, and no benefits on other decisions or later
reuse are included. If $\lambda c_*\ge K$, stopping weakly dominates
every such continuation in original objective plus computation cost.
\end{proposition}

\begin{proof}
Let $J$ be the incumbent's original cost and $J'$ the final cost of
any feasible continuation. Feasibility gives $J'\ge V_0$, so
$J-J'\le J-V_0\le K$. The continuation's incremental computation
cost is at least $\lambda c_*$, which is no smaller than that maximum
possible saving. Its total cost cannot be lower than stopping.
Sunk computation cancels from the comparison. The conclusion fails to
cover reuse benefits by construction; those require their own economic
account and cannot be added to a single-decision calculation afterwards.
\end{proof}
